\section{Applications rationnelles}
\label{section:I.7-fr}

\subsection{Applications rationnelles et fonctions rationnelles.}
\label{subsection:I.7.1-fr}

\begin{env}[7.1.1]
\label{I.7.1.1-fr}
Soient $X$, $Y$ deux préschémas, $U$, $V$, deux ouverts denses dans $X$, $f$
(resp. $g$) un morphisme de $U$ (resp. $V$) dans $Y$~; nous dirons que $f$ et
$g$ sont
\emph{équivalents} s'ils coïncident dans un ouvert dense dans $U\cap V$.
Comme une intersection finie d'ouverts denses dans $X$ est un ouvert dense
dans $X$, il est clair que cette relation est une
\emph{relation d'équivalence}.
\end{env}

\begin{definition}[7.1.2]
\label{I.7.1.2-fr}
Étant donnés deux préschémas $X$, $Y$, on appelle
\emph{application rationnelle de $X$ dans $Y$} une classe d'équivalence de
morphismes de parties ouvertes denses de $X$ dans $Y$, pour la relation
définie dans (\hyperref[I.7.1.1-fr]{7.1.1}). Si $X$ et $Y$ sont des
$S$-préschémas, on dit qu'une telle classe est une
\emph{$S$-application rationnelle} s'il existe un représentant de cette classe
qui est un $S$-morphisme. On appelle
\emph{$S$-section rationnelle de $X$} toute $S$-application rationnelle de $S$
dans $X$. On appelle
\emph{fonction rationnelle sur un préschéma $X$} toute $X$-section rationnelle
du $X$-préschéma $X\otimes_{\mathbf{Z}}\mathbf{Z}[T]$ (où $T$ est une
indéterminée).

Par abus de langage, lorsqu'il ne sera question que de $S$-préschémas, on dira
«~application rationnelle~» au lieu de «~$S$-application rationnelle~» si
aucune confusion ne peut en résulter.

Soient $f$ une application rationnelle de $X$ dans $Y$, $U$ un ouvert de
$X$~; si $f_1$, $f_2$ sont des morphismes appartenant à la classe $f$, définis
respectivement dans des ouverts denses $V$, $W$ de $X$, les restrictions
$f_1|(U\cap V)$, $f_2|(U\cap W)$ coïncident dans $U\cap V\cap W$ qui est dense
dans $U$~; la classe de morphismes $f$ définit donc une application rationnelle
de $U$ dans $Y$, appelée
\emph{restriction de $f$ à $U$} et notée $f|U$.

Si, à tout $S$-morphisme $f:X\to Y$, on fait correspondre la $S$-application
rationnelle
\oldpage[I]{156}
à laquelle appartient $f$, on définit une application canonique de
$\operatorname{Hom}_S(X,Y)$ dans l'ensemble des $S$-applications rationnelles
de $X$ dans $Y$. On désigne par $\Gamma_{\mathrm{rat}}(X/Y)$ l'ensemble des
$Y$-sections rationnelles de $X$, et on a donc une application canonique
$\Gamma(X/Y)\to\Gamma_{\mathrm{rat}}(X/Y)$. Il est clair en outre que si $X$ et
$Y$ sont deux $S$-préschémas, l'ensemble des $S$-applications rationnelles de
$X$ dans $Y$ s'identifie canoniquement à
$\Gamma_{\mathrm{rat}}((X\times_S Y)/X)$
(\hyperref[I.3.3.14-fr]{3.3.14}).
\end{definition}

\begin{env}[7.1.3]
\label{I.7.1.3-fr}
Il résulte aussitôt de (\hyperref[I.7.1.2-fr]{7.1.2}) et de
(\hyperref[I.3.3.14-fr]{3.3.14}) que les fonctions rationnelles sur $X$
s'identifient canoniquement aux
\emph{classes d'équivalence des sections du faisceau structural
$\mathscr{O}_X$} au-dessus d'ouverts partout denses de $X$, deux telles
sections étant équivalentes si elles coïncident dans un ouvert partout dense
contenu dans l'intersection de leurs ensembles de définition. Il en résulte en
particulier que les fonctions rationnelles sur $X$ forment un
\emph{anneau} $R(X)$.
\end{env}

\begin{env}[7.1.4]
\label{I.7.1.4-fr}
Lorsque $X$ est un préschéma \emph{irréductible}, tout ouvert non vide est dense
dans $X$~; on peut encore dire que les ouverts non vides de $X$ sont les
\emph{voisinages ouverts du point générique $x$} de $X$. Dire que deux
morphismes de parties ouvertes non vides de $X$ dans $Y$ sont équivalents
signifie donc dans ce cas qu'ils ont \emph{même germe} au point $x$. Autrement
dit, les applications rationnelles (resp. $S$-applications rationnelles)
$X\to Y$ s'identifient aux \emph{germes de morphismes} (resp. de
$S$-morphismes) de parties ouvertes non vides de $X$ dans $Y$ au point
générique $x$ de $X$. En particulier~:
\end{env}

\begin{proposition}[7.1.5]
\label{I.7.1.5-fr}
Si $X$ est un préschéma irréductible, l'anneau $R(X)$ des fonctions rationnelles
sur $X$ s'identifie canoniquement à l'anneau local $\mathscr{O}_x$ du point
générique $x$ de $X$. C'est un anneau local de dimension 0, et par suite un
anneau local artinien lorsque $X$ est noethérien~; c'est un corps lorsque $X$
est intègre, et il s'identifie au corps des fractions de $A(X)$ lorsqu'en outre
$X$ est un schéma affine.
\end{proposition}

Vu ce qui précède et l'identification des fonctions rationnelles aux sections
de $\mathscr{O}_X$ au-dessus d'un ouvert partout dense, la première assertion
n'est autre que la définition de la fibre d'un faisceau en un point. Pour les
autres assertions, on peut se borner au cas où $X$ est affine d'anneau $A$~;
alors $\mathfrak{j}_x$ est le nilradical de $A$, et $\mathscr{O}_x$ est donc de
dimension 0~; si $A$ est intègre, $\mathfrak{j}_x=(0)$, et $\mathscr{O}_x$ est
donc le corps des fractions de $A$. Enfin, si $A$ est noethérien, on sait
(\cite{I-11}, p.~127, cor.~4) que $\mathfrak{j}_x$ est nilpotent et
$\mathscr{O}_x=A_{\mathfrak{j}_x}$ artinien.

Si $X$ est \emph{intègre}, l'anneau $\mathscr{O}_z$ est intègre pour tout
$z\in X$~; tout ouvert affine $U$ contenant $z$ contient aussi $x$, et $R(U)$,
égal au corps des fractions de $A(U)$, s'identifie donc à $R(X)$~; on en
conclut que $R(X)$ s'identifie aussi au
\emph{corps des fractions de $\mathscr{O}_z$}~: l'identification canonique de
$\mathscr{O}_z$ à un sous-anneau de $R(X)$ consiste à faire correspondre à tout
germe de section $s\in\mathscr{O}_z$ l'unique fonction rationnelle sur $X$,
classe d'une section de $\mathscr{O}_X$ (nécessairement définie sur un ouvert
partout dense) ayant pour germe $s$ au point $z$.

\begin{env}[7.1.6]
\label{I.7.1.6-fr}
Supposons maintenant que $X$ ait un nombre \emph{fini} de composantes
irréductibles $X_i$ ($1\leq i\leq n$) (ce qui est le cas lorsque l'espace
sous-jacent à $X$ est \emph{noethérien})~; soit $X'_i$ l'ouvert de $X$,
complémentaire par rapport à $X_i$ de la réunion des $X_j\cap X_i$ pour
$j\neq i$~; $X'_i$ est irréductible, son point générique $x_i$ est le point
générique de $X_i$, et les $X'_i$ sont deux à deux disjoints, leur réunion
étant dense dans $X$ (\hyperref[0.2.1.6-fr]{0, 2.1.6}). Pour
\oldpage[I]{157}
tout ouvert partout dense $U$ de $X$, $U_i=U\cap X'_i$ est un ouvert non vide
dense dans $X'_i$, les $U_i$ étant deux à deux sans point commun, donc
$U'=\bigcup_i U_i$ est dense dans $X$. Se donner un morphisme de $U'$ dans $Y$
revient à se donner (arbitrairement) un morphisme de chacun des $U_i$ dans $Y$.
Donc~:
\end{env}

\begin{proposition}[7.1.7]
\label{I.7.1.7-fr}
Soient $X$, $Y$ deux préschémas (resp. $S$-préschémas) tels que $X$ ait un
nombre fini de composantes irréductibles $X_i$, de points génériques $x_i$
($1\leq i\leq n$). Si $R_i$ est l'ensemble des germes de morphismes (resp.
$S$-morphismes) de parties ouvertes de $X$ dans $Y$ au point $x_i$, l'ensemble
des applications rationnelles (resp. $S$-applications rationnelles) de $X$
dans $Y$ s'identifie au produit des $R_i$ ($1\leq i\leq n$).
\end{proposition}

\begin{corollary}[7.1.8]
\label{I.7.1.8-fr}
Soit $X$ un préschéma noethérien. L'anneau des fonctions rationnelles sur $X$
est un anneau artinien, dont les composants locaux sont les anneaux
$\mathscr{O}_{x_i}$ des points génériques $x_i$ des composantes irréductibles
de $X$.
\end{corollary}

\begin{corollary}[7.1.9]
\label{I.7.1.9-fr}
Soient $A$ un anneau noethérien, et soit $X=\operatorname{Spec}(A)$. Si $Q$ est
le complémentaire de la réunion des idéaux premiers minimaux de $A$, l'anneau
des fonctions rationnelles sur $X$ s'identifie canoniquement à l'anneau de
fractions $Q^{-1}A$.
\end{corollary}

Cela résultera du lemme suivant~:

\begin{lemma}[7.1.9.1]
\label{I.7.1.9.1-fr}
Pour qu'un élément $f\in A$ soit tel que $D(f)$ soit partout dense dans $X$, il
faut et il suffit que $f\in Q$~; tout ouvert dense dans $X$ contient un ouvert
de la forme $D(f)$, où $f\in Q$.
\end{lemma}

En effet, supposons ce lemme démontré~; l'anneau des sections
$\Gamma(D(f),\mathscr{O}_X)$ s'identifiant à $A_f$
(\hyperref[I.1.3.6-fr]{1.3.6} et \hyperref[I.1.3.7-fr]{1.3.7}), il résulte du
fait que les $D(f)$ avec $f\in Q$ forment un ensemble cofinal dans l'ensemble
ordonné (pour $\supset$) des ouverts denses dans $X$, et de la déf.
(\hyperref[I.7.1.1-fr]{7.1.1}) que l'anneau des fonctions rationnelles sur $X$
s'identifie à la limite inductive des $A_f$ pour $f\in Q$ (pour la relation de
préordre «~$g$ est multiple de $f$~»), c'est-à-dire à $Q^{-1}A$
(\hyperref[0.1.4.5-fr]{0, 1.4.5}).

Pour démontrer (\hyperref[I.7.1.9.1-fr]{7.1.9.1}), désignons encore par $X_i$
les composantes irréductibles de $X$ ($1\leq i\leq n$)~; si $D(f)$ est dense
dans $X$, $D(f)\cap X_i\neq\varnothing$ pour $1\leq i\leq n$ et
réciproquement~; mais cela signifie que $f\notin\mathfrak{p}_i$ pour
$1\leq i\leq n$, en posant $\mathfrak{p}_i=\mathfrak{j}(X_i)$, et comme les
$\mathfrak{p}_i$ sont les idéaux premiers minimaux de $A$
(\hyperref[I.1.1.14-fr]{1.1.14}), les relations
$f\notin\mathfrak{p}_i$ ($1\leq i\leq n$) équivalent à $f\in Q$, d'où la
première assertion du lemme. D'autre part, si $U$ est un ouvert dense dans $X$,
le complémentaire de $U$ est un ensemble de la forme $V(\mathfrak{a})$, où
$\mathfrak{a}$ est un idéal qui n'est contenu dans aucun des $\mathfrak{p}_i$~;
il n'est donc pas contenu dans leur réunion (\cite{I-10}, p.~13), et il existe
donc un $f\in\mathfrak{a}$ appartenant à $Q$~; d'où $D(f)\subset U$, ce qui
achève la démonstration.

\begin{env}[7.1.10]
\label{I.7.1.10-fr}
Supposons de nouveau $X$ irréductible, de point générique $x$. Comme tout
ouvert non vide $U$ de $X$ contient $x$, et par suite contient aussi tout
$z\in X$ tel que $x\in\overline{\{z\}}$, tout morphisme $U\to Y$ peut se
composer avec le morphisme canonique $\operatorname{Spec}(\mathscr{O}_x)\to X$
(\hyperref[I.2.4.1-fr]{2.4.1})~; et deux morphismes dans $Y$ de deux parties
ouvertes non vides de $X$, qui coïncident dans une partie ouverte non vide de
$X$ donnent par composition le même morphisme
$\operatorname{Spec}(\mathscr{O}_x)\to Y$. Autrement dit, à toute application
rationnelle de $X$ dans $Y$ correspond ainsi un morphisme
$\operatorname{Spec}(\mathscr{O}_x)\to Y$ bien déterminé.
\end{env}

\oldpage[I]{158}
\begin{proposition}[7.1.11]
\label{I.7.1.11-fr}
Soient $X$, $Y$ deux $S$-préschémas~; on suppose $X$ irréductible, de point
générique $x$, et $Y$ de type fini sur $S$. Deux $S$-applications rationnelles
de $X$ dans $Y$, auxquelles correspond le même $S$-morphisme
$\operatorname{Spec}(\mathscr{O}_x)\to Y$ sont identiques. Si on suppose de
plus $S$ localement noethérien, tout $S$-morphisme de
$\operatorname{Spec}(\mathscr{O}_x)$ dans $Y$ correspond à une $S$-application
rationnelle (et une seule) de $X$ dans $Y$.
\end{proposition}

Compte tenu de ce que tout ouvert non vide de $X$ est partout dense, cela
résulte aussitôt de (\hyperref[I.6.5.1-fr]{6.5.1}).

\begin{corollary}[7.1.12]
\label{I.7.1.12-fr}
On suppose $S$ localement noethérien, et les autres hypothèses de
(\hyperref[I.7.1.11-fr]{7.1.11}) satisfaites. Les $S$-applications rationnelles
de $X$ dans $Y$ s'identifient alors aux points du $S$-préschéma $Y$, à valeurs
dans le $S$-préschéma $\operatorname{Spec}(\mathscr{O}_x)$.
\end{corollary}

Cela n'est autre que (\hyperref[I.7.1.11-fr]{7.1.11}), avec la terminologie
introduite dans (\hyperref[I.3.4.1-fr]{3.4.1}).

\begin{corollary}[7.1.13]
\label{I.7.1.13-fr}
On suppose remplies les conditions de (\hyperref[I.7.1.12-fr]{7.1.12}). Soit
$s$ l'image de $x$ dans $S$. La donnée d'une $S$-application rationnelle de
$X$ dans $Y$ équivaut à la donnée d'un point $y$ de $Y$ au-dessus de $s$, et
d'un $\mathscr{O}_s$-homomorphisme local
$\mathscr{O}_y\to\mathscr{O}_x=R(X)$.
\end{corollary}

Cela résulte de (\hyperref[I.7.1.11-fr]{7.1.11}) et de
(\hyperref[I.2.4.4-fr]{2.4.4}).

En particulier~:

\begin{corollary}[7.1.14]
\label{I.7.1.14-fr}
Sous les conditions de (\hyperref[I.7.1.12-fr]{7.1.12}), les $S$-applications
rationnelles de $X$ dans $Y$ ne dépendent (pour $Y$ donné) que du
$S$-préschéma $\operatorname{Spec}(\mathscr{O}_x)$ et en particulier restent
les mêmes lorsqu'on remplace $X$ par
$\operatorname{Spec}(\mathscr{O}_z)$, pour tout $z\in X$.
\end{corollary}

En effet, comme $z\in\overline{\{x\}}$, $x$ est le point générique de
$Z=\operatorname{Spec}(\mathscr{O}_z)$, et
$\mathscr{O}_{X,x}=\mathscr{O}_{Z,x}$.

Lorsque $X$ est intègre, $R(X)=\mathscr{O}_x=k(x)$ est un corps
(\hyperref[I.7.1.5-fr]{7.1.5})~; les corollaires précédents se spécialisent
alors en~:

\begin{corollary}[7.1.15]
\label{I.7.1.15-fr}
On suppose vérifiées les conditions de (\hyperref[I.7.1.12-fr]{7.1.12}) et en
outre que $X$ est intègre. Soit $s$ l'image de $x$ dans $S$. Alors les
$S$-applications rationnelles de $X$ dans $Y$ s'identifient aux points
géométriques de $Y\otimes_S k(s)$ à valeurs dans l'extension $R(X)$ de $k(s)$,
autrement dit chacune d'elles équivaut à la donnée d'un point $y\in Y$
au-dessus de $s$ et d'un $k(s)$-monomorphisme de $k(y)$ dans
$k(x)=R(X)$.
\end{corollary}

Les points de $Y$ au-dessus de $s$ s'identifient en effet à ceux de
$Y\otimes_S k(s)$ (\hyperref[I.3.6.3-fr]{3.6.3}) et les
$\mathscr{O}_s$-homomorphismes locaux $\mathscr{O}_y\to R(X)$ aux
$k(s)$-monomorphismes $k(y)\to R(X)$.

Plus particulièrement~:

\begin{corollary}[7.1.16]
\label{I.7.1.16-fr}
Soient $k$ un corps, $X$, $Y$ deux préschémas algébriques
(\hyperref[I.6.4.1-fr]{6.4.1}) sur $k$~; on suppose en outre $X$ intègre. Alors
les $k$-applications rationnelles de $X$ dans $Y$ s'identifient aux points
géométriques de $Y$ à valeurs dans l'extension $R(X)$ de $k$
(\hyperref[I.3.4.4-fr]{3.4.4}).
\end{corollary}

\subsection{Domaine de définition d'une application rationnelle.}
\label{subsection:I.7.2-fr}

\begin{env}[7.2.1]
\label{I.7.2.1-fr}
Soient $X$, $Y$ deux préschémas, $f$ une application rationnelle de $X$ dans
$Y$. On dit que $f$ est \emph{définie en un point $x\in X$} s'il existe un
ensemble ouvert partout dense $U$ contenant $x$ et un morphisme $U\to Y$
appartenant à la classe d'équivalence $f$. L'ensemble des points $x\in X$ où
$f$ est définie est appelé le \emph{domaine de définition de $f$}~; il est
clair que c'est un ouvert partout dense dans $X$.
\end{env}

\oldpage[I]{159}
\begin{proposition}[7.2.2]
\label{I.7.2.2-fr}
Soient $X$, $Y$ deux $S$-préschémas, tels que $X$ soit réduit et $Y$ séparé
sur $S$. Soit $f$ une $S$-application rationnelle de $X$ dans $Y$, $U_0$ son
domaine de définition. Il existe alors un $S$-morphisme et un seul $U_0\to Y$
appartenant à la classe $f$.
\end{proposition}

Comme pour tout morphisme $U\to Y$ appartenant à la classe $f$, on a
nécessairement $U\subset U_0$, il est clair que la proposition sera
conséquence du

\begin{lemma}[7.2.2.1]
\label{I.7.2.2.1-fr}
Sous les hypothèses de (\hyperref[I.7.2.2-fr]{7.2.2}), soient $U_1$, $U_2$
deux ouverts partout denses de $X$, $f_i:U_i\to Y$ ($i=1,2$) deux
$S$-morphismes tels qu'il existe un ouvert $V\subset U_1\cap U_2$, dense dans
$X$ et dans lequel $f_1$ et $f_2$ coïncident. Alors $f_1$ et $f_2$ coïncident
dans $U_1\cap U_2$.
\end{lemma}

On peut évidemment se borner au cas où $X=U_1=U_2$. Comme $X$ (donc $V$)
est réduit, $X$ est le plus petit sous-préschéma fermé de $X$ majorant $V$
(\hyperref[I.5.2.2-fr]{5.2.2}). Soit
$g=(f_1,f_2)_S:X\to Y\times_S Y$~; comme par hypothèse la diagonale
$T=\Delta_Y(Y)$ est un sous-préschéma fermé de $Y\times_S Y$,
$Z=g^{-1}(T)$ est un sous-préschéma fermé de $X$
(\hyperref[I.4.4.1-fr]{4.4.1}). Si $h:V\to Y$ est la restriction commune de
$f_1$ et $f_2$ à $V$, la restriction de $g$ à $V$ est $g'=(h,h)_S$, qui se
factorise en $g'=\Delta_Y\circ h$~; comme $\Delta_Y^{-1}(T)=Y$, on a
$g'^{-1}(T)=V$, et par suite $Z$ est un sous-préschéma fermé de $X$ induisant
$V$, donc majorant $V$, ce qui entraîne $Z=X$. De la relation
$g^{-1}(T)=X$, on déduit (\hyperref[I.4.4.1-fr]{4.4.1}) que $g$ se factorise
en $\Delta_Y\circ f$, où $f$ est un morphisme $X\to Y$, ce qui entraîne par
définition du morphisme diagonal que $f_1=f_2=f$.

Il est clair que le morphisme $U_0\to Y$ défini dans
(\hyperref[I.7.2.2-fr]{7.2.2}) est l'unique morphisme de la classe $f$ qui
\emph{ne puisse être prolongé} à un morphisme d'une partie ouverte de $X$
contenant strictement $U_0$. \emph{Sous les hypothèses de
(\hyperref[I.7.2.2-fr]{7.2.2})}, on peut donc \emph{identifier} les
applications rationnelles de $X$ dans $Y$ aux morphismes \emph{non
prolongeables} (à des ouverts strictement plus grands) d'ouverts partout
denses de $X$ dans $Y$. Avec cette identification, la prop.
(\hyperref[I.7.2.2-fr]{7.2.2}) entraîne~:

\begin{corollary}[7.2.3]
\label{I.7.2.3-fr}
Les hypothèses sur $X$ et $Y$ étant celles de
(\hyperref[I.7.2.2-fr]{7.2.2}), soit $U$ un ouvert partout dense de $X$. Il
existe une correspondance biunivoque canonique entre les $S$-morphismes de
$U$ dans $Y$ et les $S$-applications rationnelles de $X$ dans $Y$ définies
en tous les points de $U$.
\end{corollary}

En vertu de (\hyperref[I.7.2.2-fr]{7.2.2}), pour tout $S$-morphisme $f$ de
$U$ dans $Y$, il existe en effet une $S$-application rationnelle et une seule
$\overline f$ de $X$ dans $Y$ qui prolonge $f$.

\begin{corollary}[7.2.4]
\label{I.7.2.4-fr}
Soient $S$ un schéma, $X$ un $S$-préschéma réduit, $Y$ un $S$-schéma,
$f:U\to Y$ un $S$-morphisme d'un ouvert dense $U$ de $X$ dans $Y$. Si
$\overline f$ est la $\mathbf{Z}$-application rationnelle de $X$ dans $Y$ qui
prolonge $f$, $\overline f$ est un $S$-morphisme (et est par suite la
$S$-application rationnelle de $X$ dans $Y$ prolongeant $f$).
\end{corollary}

En effet, si $\varphi:X\to S$, $\psi:Y\to S$ sont les morphismes structuraux,
$U_0$ le domaine de définition de $\overline f$, $j$ l'injection $U_0\to X$,
il suffit de prouver que $\psi\circ\overline f=\varphi\circ j$, ce qui résulte
aussitôt de (\hyperref[I.7.2.2.1-fr]{7.2.2.1}), puisque $f$ est un
$S$-morphisme.

\begin{corollary}[7.2.5]
\label{I.7.2.5-fr}
Soient $X$, $Y$ deux $S$-préschémas~; on suppose $X$ réduit, $X$ et $Y$
séparés sur $S$. Soient $p:Y\to X$ un $S$-morphisme (faisant de $Y$ un
$X$-préschéma), $U$ un ouvert partout dense de $X$, $f$ une $U$-section de
$Y$~; alors l'application rationnelle $\overline f$ de $X$ dans $Y$
prolongeant $f$ est une $X$-section rationnelle de $Y$.
\end{corollary}

\oldpage[I]{160}
Il faut prouver que $p\circ\overline f$ est l'identité dans le domaine de
définition de $\overline f$~; puisque $X$ est séparé sur $S$, cela résulte
encore de (\hyperref[I.7.2.2.1-fr]{7.2.2.1}).

\begin{corollary}[7.2.6]
\label{I.7.2.6-fr}
Soient $X$ un préschéma réduit, $U$ un ouvert partout dense de $X$. Il y a
correspondance biunivoque canonique entre les sections de $\mathscr{O}_X$ au-dessus
de $U$ et les fonctions rationnelles $f$ sur $X$ définies en tout point de $U$.
\end{corollary}

Compte tenu de (\hyperref[I.7.2.3-fr]{7.2.3}),
(\hyperref[I.7.1.2-fr]{7.1.2}) et (\hyperref[I.7.1.3-fr]{7.1.3}), il suffit de
remarquer que le $X$-préschéma $X\otimes_{\mathbf Z}\mathbf Z[T]$ est séparé
au-dessus de $X$ (\hyperref[I.5.5.1-fr]{5.5.1, (iv)}).

\begin{corollary}[7.2.7]
\label{I.7.2.7-fr}
Soient $Y$ un préschéma réduit, $f:X\to Y$ un morphisme séparé, $U$ un ouvert
partout dense de $Y$, $g:U\to f^{-1}(U)$ une $U$-section de $f^{-1}(U)$, $Z$
le sous-préschéma réduit de $X$ ayant $\overline{g(U)}$ pour espace sous-jacent
(\hyperref[I.5.2.1-fr]{5.2.1}). Pour que $g$ soit restriction d'une
$Y$-section de $X$ (autrement dit (\hyperref[I.7.2.5-fr]{7.2.5}) pour que
l'application rationnelle de $Y$ dans $X$ prolongeant $g$ soit partout
définie), il faut et il suffit que la restriction de $f$ à $Z$ soit un
isomorphisme de $Z$ sur $Y$.
\end{corollary}

La restriction de $f$ à $f^{-1}(U)$ est un morphisme séparé
(\hyperref[I.5.5.1-fr]{5.5.1, (i)}), donc $g$ est une immersion fermée
(\hyperref[I.5.4.6-fr]{5.4.6}), et par suite
$g(U)=Z\cap f^{-1}(U)$ et le sous-préschéma induit par $Z$ sur l'ouvert
$g(U)$ de $Z$ est identique au sous-préschéma fermé de $f^{-1}(U)$ associé à
$g$ (\hyperref[I.5.2.1-fr]{5.2.1}). Il est clair alors que la condition de
l'énoncé est suffisante, car si elle est remplie et si $f_Z:Z\to Y$ est la
restriction de $f$ à $Z$ et $\overline g:Y\to Z$ l'isomorphisme réciproque,
$\overline g$ prolonge $g$. Inversement, si $g$ est restriction à $U$ d'une
$Y$-section $h$ de $X$, $h$ est une immersion fermée
(\hyperref[I.5.4.6-fr]{5.4.6}), donc $h(Y)$ est fermé, et comme il est contenu
dans $Z$, il est égal à $Z$, et il résulte de
(\hyperref[I.5.2.1-fr]{5.2.1}) que $h$ est nécessairement un isomorphisme de
$Y$ sur le sous-préschéma fermé $Z$ de $X$.

\begin{env}[7.2.8]
\label{I.7.2.8-fr}
Soient $X$, $Y$ deux $S$-préschémas, $X$ étant supposé réduit et $Y$ séparé
sur $S$. Soit $f$ une $S$-application rationnelle de $X$ dans $Y$, et soit
$x$ un point de $X$~; on peut composer $f$ avec le $S$-morphisme canonique
$\operatorname{Spec}(\mathscr{O}_x)\to X$
(\hyperref[I.2.4.1-fr]{2.4.1}) pourvu que la trace sur
$\operatorname{Spec}(\mathscr{O}_x)$ du domaine de définition de $f$ soit
dense dans $\operatorname{Spec}(\mathscr{O}_x)$ (identifié à l'ensemble des
$z\in X$ tels que $x\in\overline{\{z\}}$
(\hyperref[I.2.4.2-fr]{2.4.2})). Ceci aura lieu dans les cas suivants~:
\begin{enumerate}
  \item[1\textsuperscript{o}] $X$ est irréductible (donc intègre), car alors
    le point générique $\xi$ de $X$ est le point générique de
    $\operatorname{Spec}(\mathscr{O}_x)$~; comme le domaine de définition $U$
    de $f$ contient $\xi$, $U\cap\operatorname{Spec}(\mathscr{O}_x)$ contient
    $\xi$, donc est dense dans $\operatorname{Spec}(\mathscr{O}_x)$.
  \item[2\textsuperscript{o}] $X$ est localement noethérien~; notre assertion
    résulte en effet alors du
\end{enumerate}
\end{env}

\begin{lemma}[7.2.8.1]
\label{I.7.2.8.1-fr}
Soient $X$ un préschéma dont l'espace sous-jacent est localement noethérien,
$x$ un point de $X$. Les composantes irréductibles de
$\operatorname{Spec}(\mathscr{O}_x)$ sont les traces sur
$\operatorname{Spec}(\mathscr{O}_x)$ des composantes irréductibles de $X$
contenant $x$. Pour qu'un ouvert $U\subset X$ soit tel que
$U\cap\operatorname{Spec}(\mathscr{O}_x)$ soit dense dans
$\operatorname{Spec}(\mathscr{O}_x)$, il faut et il suffit qu'il rencontre les
composantes irréductibles de $X$ contenant $x$ (ce qui a lieu en particulier
si $U$ est dense dans $X$).
\end{lemma}

La seconde assertion résulte évidemment de la première, et il suffit donc de
démontrer celle-ci. Comme $\operatorname{Spec}(\mathscr{O}_x)$ est contenu
dans tout ouvert affine $U$ contenant $x$, et que les composantes
irréductibles de $U$ contenant $x$ sont les traces sur $U$ des composantes
irréductibles de $X$ contenant $x$
(\hyperref[0.2.1.6-fr]{0, 2.1.6}), on peut supposer $X$ affine d'anneau $A$.
Comme les idéaux premiers de $A_x$ correspondent biunivoquement aux idéaux
premiers
\oldpage[I]{161}
de $A$ contenus dans $\mathfrak{j}_x$
(\hyperref[0.1.2.6-fr]{0, 1.2.6}), les idéaux premiers minimaux de $A_x$
correspondent aux idéaux premiers minimaux de $A$ contenus dans
$\mathfrak{j}_x$, d'où le lemme (\hyperref[I.1.1.14-fr]{1.1.14}).

Cela étant, supposons que l'on soit dans l'un des deux cas précités. Si $U$
est le domaine de définition de la $S$-application rationnelle $f$, désignons
par $f'$ l'application rationnelle de $\operatorname{Spec}(\mathscr{O}_x)$
dans $Y$ qui coïncide (compte tenu de (\hyperref[I.2.4.2-fr]{2.4.2})) avec
$f$ dans $U\cap\operatorname{Spec}(\mathscr{O}_x)$~; nous dirons que cette
application rationnelle est \emph{induite par $f$}.

\begin{proposition}[7.2.9]
\label{I.7.2.9-fr}
Soient $S$ un préschéma localement noethérien, $X$ un $S$-préschéma réduit,
$Y$ un $S$-schéma de type fini. On suppose en outre $X$ irréductible ou
localement noethérien. Soient alors $f$ une $S$-application rationnelle de
$X$ dans $Y$, $x$ un point de $X$. Pour que $f$ soit définie au point $x$,
il faut et il suffit que l'application rationnelle $f'$ de
$\operatorname{Spec}(\mathscr{O}_x)$ dans $Y$, induite par $f$
(\hyperref[I.7.2.8-fr]{7.2.8}), soit un morphisme.
\end{proposition}

La condition étant évidemment nécessaire (puisque
$\operatorname{Spec}(\mathscr{O}_x)$ est contenu dans tout ouvert contenant
$x$), prouvons qu'elle est suffisante. En vertu de
(\hyperref[I.6.5.1-fr]{6.5.1}), il existe un voisinage ouvert $U$ de $x$ dans
$X$ et un $S$-morphisme $g$ de $U$ dans $Y$, induisant $f'$ sur
$\operatorname{Spec}(\mathscr{O}_x)$. Si $X$ est irréductible, $U$ est dense
dans $X$, et en vertu de (\hyperref[I.7.2.3-fr]{7.2.3}) on peut supposer que
$g$ est une $S$-application rationnelle. En outre, le point générique de $X$
appartient à $\operatorname{Spec}(\mathscr{O}_x)$ et au domaine de définition
de $f$, donc $f$ et $g$ coïncident en ce point, et par suite dans un ensemble
ouvert non vide de $X$ (\hyperref[I.6.5.1-fr]{6.5.1}). Mais comme $f$ et $g$
sont des $S$-applications rationnelles, elles sont identiques
(\hyperref[I.7.2.3-fr]{7.2.3}), donc $f$ est définie en $x$.

Si maintenant on suppose $X$ localement noethérien, on peut supposer $U$
noethérien~; il n'y a alors qu'un nombre fini de composantes irréductibles
$X_i$ de $X$ contenant $x$ (\hyperref[I.7.2.8.1-fr]{7.2.8.1}), et on peut
supposer que ce sont les seules rencontrant $U$, en remplaçant au besoin $U$
par un ouvert plus petit (puisqu'il n'y a qu'un nombre fini de composantes
irréductibles de $X$ rencontrant $U$, $U$ étant noethérien). On voit alors
comme ci-dessus que $f$ et $g$ coïncident dans un ouvert non vide de chacun
des $X_i$. Tenant compte du fait que chacun des $X_i$ est contenu dans
$\overline U$, considérons alors le morphisme $f_1$, défini dans un ouvert
dense de $U\cup(X-\overline U)$, égal à $g$ dans $U$ et à $f$ dans
l'intersection de $X-\overline U$ et du domaine de définition de $f$. Comme
$U\cup(X-\overline U)$ est dense dans $X$, $f_1$ et $f$ coïncident dans un
ouvert dense de $X$, et comme $f$ est une application rationnelle, $f$ est
une extension de $f_1$ (\hyperref[I.7.2.3-fr]{7.2.3}), donc est définie au
point $x$.

\subsection{Faisceau des fonctions rationnelles.}
\label{subsection:I.7.3-fr}

\begin{env}[7.3.1]
\label{I.7.3.1-fr}
Soit $X$ un préschéma. Pour tout ouvert $U\subset X$, désignons par $R(U)$
l'anneau des fonctions rationnelles sur $U$
(\hyperref[I.7.1.3-fr]{7.1.3})~; c'est une
$\Gamma(U,\mathscr{O}_X)$-algèbre. En outre, si $V\subset U$ est un second
ouvert de $X$, toute section de $\mathscr{O}_X$ au-dessus d'une partie
ouverte partout dense de $U$ donne par restriction à $V$ une section
au-dessus d'une partie ouverte partout dense de $V$, et si deux sections
coïncident au-dessus d'une partie ouverte partout dense de $U$, leurs
restrictions à $V$ coïncident au-dessus d'une partie ouverte partout dense de
$V$. On définit donc ainsi un di-homomorphisme d'algèbres
$\rho_{V,U}:R(U)\to R(V)$, et il est clair que si $U\supset V\supset W$ sont
trois ouverts de $X$, on a
$\rho_{W,U}=\rho_{W,V}\circ\rho_{V,U}$~; les $R(U)$ définissent donc un
\emph{préfaisceau} d'algèbres sur $X$.
\end{env}

\oldpage[I]{162}
\begin{definition}[7.3.2]
\label{I.7.3.2-fr}
On appelle \emph{faisceau des fonctions rationnelles} sur un préschéma $X$
et on désigne par $\mathscr{R}(X)$ la $\mathscr{O}_X$-Algèbre associée au
préfaisceau formé des $R(U)$.
\end{definition}

Pour tout préschéma $X$ et tout ouvert $U\subset X$, il est clair que le
faisceau induit $\mathscr{R}(X)|U$ n'est autre que $\mathscr{R}(U)$.

\begin{proposition}[7.3.3]
\label{I.7.3.3-fr}
Soit $X$ un préschéma tel que la famille $(X_\lambda)$ de ses composantes
irréductibles soit localement finie (ce qui est en particulier le cas lorsque
l'espace sous-jacent à $X$ est localement noethérien). Alors le
$\mathscr{O}_X$-Module $\mathscr{R}(X)$ est quasi-cohérent, et pour tout
ouvert $U$ de $X$ ne rencontrant qu'un nombre fini de composantes $X_\lambda$,
$R(U)$ est égal à $\Gamma(U,\mathscr{R}(X))$ et s'identifie canoniquement au
composé direct des anneaux locaux des points génériques des $X_\lambda$ telles
que $U\cap X_\lambda\neq\varnothing$.
\end{proposition}

On peut évidemment se limiter au cas où $X$ n'a qu'un nombre fini de
composantes irréductibles $X_i$, de points génériques $x_i$
($1\leq i\leq n$). Le fait que $R(U)$ est canoniquement identifié au composé
direct des $\mathscr{O}_{x_i}=R(X_i)$ tels que
$U\cap X_i\neq\varnothing$ résulte alors de
(\hyperref[I.7.1.7-fr]{7.1.7}). Montrons en outre que le préfaisceau
$U\to R(U)$ vérifie les axiomes des faisceaux, ce qui prouvera que
$R(U)=\Gamma(U,\mathscr{R}(X))$. En effet, il vérifie (F 1) d'après ce qui
précède. Pour voir qu'il satisfait à (F 2), considérons un recouvrement d'un
ouvert $U$ de $X$ par des ouverts $V_\alpha\subset U$~; si les
$s_\alpha\in R(V_\alpha)$ sont telles que les restrictions de $s_\alpha$ et
$s_\beta$ à $V_\alpha\cap V_\beta$ coïncident pour tout couple d'indices, on
en conclut que pour tout indice $i$ tel que $U\cap X_i\neq\varnothing$, les
composantes dans $R(X_i)$ de toutes les $s_\alpha$ telles que
$V_\alpha\cap X_i\neq\varnothing$ sont les mêmes~; désignant par $t_i$ cette
composante, il est clair que l'élément de $R(U)$ ayant les $t_i$ pour
composantes a pour restriction $s_\alpha$ à chaque $V_\alpha$. Enfin, pour
voir que $\mathscr{R}(X)$ est quasi-cohérent, on peut se limiter au cas où
$X=\operatorname{Spec}(A)$ est affine~; en prenant pour $U$ les ouverts
affines de la forme $D(f)$, où $f\in A$, il résulte de ce qui précède et de la
définition (\hyperref[I.1.3.4-fr]{1.3.4}) que l'on a
$\mathscr{R}(X)=\widetilde M$, où $M$ est somme directe des $A$-modules
$A_{x_i}$.

\begin{corollary}[7.3.4]
\label{I.7.3.4-fr}
Soit $X$ un préschéma réduit n'ayant qu'un nombre fini de composantes
irréductibles, et soient $X_i$ ($1\leq i\leq n$) les sous-préschémas fermés
réduits de $X$ ayant pour espaces sous-jacents les composantes irréductibles de
$X$ (\hyperref[I.5.2.1-fr]{5.2.1}). Si $h_i$ est l'injection canonique
$X_i\to X$, $\mathscr{R}(X)$ est alors composée directe des
$\mathscr{O}_X$-Algèbres $(h_i)_*(\mathscr{R}(X_i))$.
\end{corollary}

\begin{corollary}[7.3.5]
\label{I.7.3.5-fr}
Si $X$ est irréductible, tout $\mathscr{R}(X)$-Module quasi-cohérent
$\mathscr{F}$ est un faisceau simple.
\end{corollary}

Il suffit de montrer que tout $x\in X$ admet un voisinage $U$ tel que
$\mathscr{F}|U$ soit un faisceau simple
(\hyperref[0.3.6.2-fr]{0, 3.6.2}), autrement dit on est ramené au cas où $X$
est affine~; on peut en outre supposer que $\mathscr{F}$ est le conoyau d'un
homomorphisme
$(\mathscr{R}(X))^{(I)}\to(\mathscr{R}(X))^{(J)}$
(\hyperref[0.5.1.3-fr]{0, 5.1.3}), et tout revient à voir que
$\mathscr{R}(X)$ est un faisceau simple~; mais cela est évident puisque
$\Gamma(U,\mathscr{R}(X))=R(X)$ pour tout ouvert non vide $U$, $U$ contenant
le point générique de $X$.

\begin{corollary}[7.3.6]
\label{I.7.3.6-fr}
Si $X$ est irréductible, pour tout $\mathscr{O}_X$-Module quasi-cohérent
$\mathscr{F}$, $\mathscr{F}\otimes_{\mathscr{O}_X}\mathscr{R}(X)$ est un
faisceau simple~; si en outre $X$ est réduit (donc intègre),
$\mathscr{F}\otimes_{\mathscr{O}_X}\mathscr{R}(X)$ est isomorphe à un
faisceau de la forme $(\mathscr{R}(X))^{(I)}$.
\end{corollary}

La seconde assertion résulte de ce que $R(X)$ est alors un corps.

\begin{proposition}[7.3.7]
\label{I.7.3.7-fr}
Supposons que le préschéma $X$ soit localement intègre ou localement
\oldpage[I]{163}
noethérien. Alors $\mathscr{R}(X)$ est une $\mathscr{O}_X$-Algèbre
quasi-cohérente~; si en outre $X$ est réduit (ce qui est le cas lorsque $X$
est localement intègre), l'homomorphisme canonique
$\mathscr{O}_X\to\mathscr{R}(X)$ est injectif.
\end{proposition}

La question étant locale, la première assertion résulte de
(\hyperref[I.7.3.3-fr]{7.3.3})~; la seconde résulte aussitôt de
(\hyperref[I.7.2.3-fr]{7.2.3}).

\begin{env}[7.3.8]
\label{I.7.3.8-fr}
Soient $X$, $Y$ deux préschémas ayant chacun un nombre fini de composantes
irréductibles, et soit $f:X\to Y$ un morphisme dont la restriction à
l'ensemble des points génériques des composantes irréductibles de $X$ est une
surjection sur l'ensemble des points génériques des composantes irréductibles
de $Y$. Alors on a
\[
  f^*(\mathscr{R}(Y))=\mathscr{R}(X).
  \tag{7.3.8.1}\label{I.7.3.8.1-fr}
\]
\end{env}

En effet, on est ramené (en vertu de
(\hyperref[I.7.3.3-fr]{7.3.3})) au cas où $X$ et $Y$ sont irréductibles, de
points génériques $x$, $y$, avec $f(x)=y$~; donc
\[
  (f^*(\mathscr{R}(Y)))_x
  =\mathscr{O}_y\otimes_{\mathscr{O}_y}\mathscr{O}_x
  =\mathscr{O}_x
\]
(\hyperref[0.4.3.1-fr]{0, 4.3.1}), ce qui démontre (7.3.8.1) en vertu de
(\hyperref[I.7.3.5-fr]{7.3.5}).

\subsection{Faisceaux de torsion et faisceaux sans torsion.}
\label{subsection:I.7.4-fr}

\begin{env}[7.4.1]
\label{I.7.4.1-fr}
Soit $X$ un préschéma \emph{intègre}. Pour tout $\mathscr{O}_X$-Module
$\mathscr{F}$, l'homomorphisme canonique
$\mathscr{O}_X\to\mathscr{R}(X)$ définit par tensorisation un homomorphisme
(dit encore \emph{canonique})
\[
  \mathscr{F}\to
  \mathscr{F}\otimes_{\mathscr{O}_X}\mathscr{R}(X)
\]
qui, sur chaque fibre, n'est autre que l'homomorphisme $z\to z\otimes 1$ de
$\mathscr{F}_x$ dans
$\mathscr{F}_x\otimes_{\mathscr{O}_x}R(X)$. Le \emph{noyau $\mathscr{T}$}
de cet homomorphisme est un sous-$\mathscr{O}_X$-Module de $\mathscr{F}$,
appelé \emph{faisceau de torsion de $\mathscr{F}$}~; il est quasi-cohérent si
$\mathscr{F}$ est quasi-cohérent
(\hyperref[I.4.1.1-fr]{4.1.1} et \hyperref[I.7.3.6-fr]{7.3.6}). On dit que
$\mathscr{F}$ est \emph{sans torsion} si $\mathscr{T}=0$ et que
$\mathscr{F}$ est un \emph{faisceau de torsion} si
$\mathscr{T}=\mathscr{F}$. Pour tout $\mathscr{O}_X$-Module $\mathscr{F}$,
$\mathscr{F}/\mathscr{T}$ est sans torsion. On déduit de
(\hyperref[I.7.3.5-fr]{7.3.5}) que~:
\end{env}

\begin{proposition}[7.4.2]
\label{I.7.4.2-fr}
Si $X$ est un préschéma intègre, tout $\mathscr{O}_X$-Module quasi-cohérent
sans torsion $\mathscr{F}$ est isomorphe à un sous-faisceau $\mathscr{G}$
d'un faisceau simple de la forme $(\mathscr{R}(X))^{(I)}$, engendré (en tant
que $\mathscr{R}(X)$-Module) par $\mathscr{G}$.
\end{proposition}

Le cardinal de $I$ est appelé le \emph{rang de $\mathscr{F}$}~; pour tout
ouvert affine non vide $U$ de $X$, le rang de $\mathscr{F}$ est égal au rang
de $\Gamma(U,\mathscr{F})$ en tant que $\Gamma(U,\mathscr{O}_X)$-module,
comme on le voit aussitôt en considérant le point générique de $X$, contenu
dans $U$. En particulier~:

\begin{corollary}[7.4.3]
\label{I.7.4.3-fr}
Sur un préschéma intègre $X$, tout $\mathscr{O}_X$-Module quasi-cohérent sans
torsion et de rang $1$ (en particulier tout $\mathscr{O}_X$-Module
inversible) est isomorphe à un sous-$\mathscr{O}_X$-Module de
$\mathscr{R}(X)$, et réciproquement.
\end{corollary}

\begin{corollary}[7.4.4]
\label{I.7.4.4-fr}
Soient $X$ un préschéma intègre, $\mathscr{L}$, $\mathscr{L}'$ deux
$\mathscr{O}_X$-Modules sans torsion, $f$ (resp. $f'$) une section de
$\mathscr{L}$ (resp. $\mathscr{L}'$) au-dessus de $X$. Pour que
$f\otimes f'=0$, il faut et il suffit que l'une des sections $f$, $f'$ soit
nulle.
\end{corollary}

Soit $x$ le point générique de $X$~; on a par hypothèse
$(f\otimes f')_x=f_x\otimes f'_x=0$. Comme $\mathscr{L}_x$ et
$\mathscr{L}'_x$ s'identifient à des sous-$\mathscr{O}_x$-Modules du corps
$\mathscr{O}_x$, la relation précédente entraîne $f_x=0$ ou $f'_x=0$, et par
suite $f=0$ ou $f'=0$ puisque $\mathscr{L}$ et $\mathscr{L}'$ sont sans
torsion (\hyperref[I.7.3.5-fr]{7.3.5}).

\begin{proposition}[7.4.5]
\label{I.7.4.5-fr}
Soient $X$, $Y$ deux préschémas intègres, $f:X\to Y$ un morphisme dominant.
Pour tout $\mathscr{O}_X$-Module quasi-cohérent sans torsion $\mathscr{F}$,
$f_*(\mathscr{F})$ est un $\mathscr{O}_Y$-Module sans torsion.
\end{proposition}

\oldpage[I]{164}
\begin{proof}
Comme $f_*$ est exact à gauche
(\hyperref[0.4.2.1-fr]{0, 4.2.1}), il suffit, en vertu de
(\hyperref[I.7.4.2-fr]{7.4.2}), de prouver la proposition lorsque
$\mathscr{F}=(\mathscr{R}(X))^{(I)}$. Or, tout ouvert non vide $U$ de $Y$
contient le point générique de $Y$, donc $f^{-1}(U)$ contient le point
générique de $X$ (\hyperref[0.2.1.5-fr]{0, 2.1.5}), donc on a alors
\[
  \Gamma(U,f_*(\mathscr{F}))
  =\Gamma(f^{-1}(U),\mathscr{F})=(R(X))^{(I)}~;
\]
autrement dit, $f_*(\mathscr{F})$ est le faisceau simple de fibre
$(R(X))^{(I)}$, considéré comme $\mathscr{R}(Y)$-Module, et il est évidemment
sans torsion.
\end{proof}

\begin{proposition}[7.4.6]
\label{I.7.4.6-fr}
Soient $X$ un préschéma intègre, $x$ son point générique. Pour tout
$\mathscr{O}_X$-Module quasi-cohérent de type fini $\mathscr{F}$, les
conditions suivantes sont équivalentes~: a) $\mathscr{F}$ est un faisceau de
torsion~; b) $\mathscr{F}_x=0$~; c)
$\operatorname{Supp}(\mathscr{F})\neq X$.
\end{proposition}

En vertu de (\hyperref[I.7.3.5-fr]{7.3.5}) et de
(\hyperref[I.7.4.1-fr]{7.4.1}), les relations $\mathscr{F}_x=0$ et
$\mathscr{F}\otimes_{\mathscr{O}_X}\mathscr{R}(X)=0$ sont équivalentes, donc
a) et b) sont équivalentes~; d'autre part,
$\operatorname{Supp}(\mathscr{F})$ est fermé dans $X$
(\hyperref[0.5.2.2-fr]{0, 5.2.2}), et comme tout ouvert non vide de $X$
contient $x$, b) et c) sont équivalentes.

\begin{env}[7.4.7]
\label{I.7.4.7-fr}
On étend (par abus de langage) les définitions de
(\hyperref[I.7.4.1-fr]{7.4.1}) au cas où $X$ est un préschéma
\emph{réduit} n'ayant qu'un nombre \emph{fini} de composantes irréductibles~;
il résulte alors de (\hyperref[I.7.3.4-fr]{7.3.4}) que l'équivalence de a) et
c) dans (\hyperref[I.7.4.6-fr]{7.4.6}) est encore valable pour un tel
préschéma.
\end{env}
